% P10-4c-Nachtrag, 15.09.2026: Kontextverfügbarkeit und Nutzung sowie
% Werkzeugrückgabe und Zustandswirkung getrennt; Logger am Code geprüft.
% M25, 13.09.2026: bestehende Dokumentationsbelege in den numerischen Zitierstil überführt.
% M11, 11.09.2026: Übergang zu 6.9 auf ausgewählte technische
% Prüfungen begrenzt; Kapitelverweise als Referenzen.
% ============================================================
% §6.8 Observability und technische Verifikation — v02
% (Claude, 03.09.)
% v02 (Kollegen-Runde 1, 4 Präzisierungen + Schlussfix):
% ① Artefaktarten konditional (nicht jeder Lauf hat Gates/Repair)
% ② „protokollierte Ausführungsschritte" statt „Lauf
%   rekonstruierbar" (A1-Grenze) ③ Messdraht auf die gemeinsame
%   Pipeline begrenzt (Steward-Dialog nicht verifiziert) +
%   Reasoning als „explizit ausgegebener Text" ④ OTel-Absatz ohne
%   Framework-Abwesenheits-Behauptung (positive Eigenanteil-Form)
% ⑤ Schluss epistemisch exakt. Kollegen-Urteil danach:
%   SEMANTISCH FROZEN.
% Grundlage: kapitel-6-gesamtvertrag.md ★FROZEN (§4 Zeile 6.8 mit
% Leitfrage v03: „Wie wird eine tatsächliche Ausführung technisch
% beobachtbar und nachweisbar?"; Kernaussage #7; Label-Regel
% sec:realisierung-observability — 3.2-/Anhang-Verweise beim
% Label-Pass auf \ref) + V6-Check #13 (UseOpenTelemetry/GenAI-
% Konventionen BESTÄTIGT) + Code-Fundstellen aus den 6.3-Checks.
% ABGRENZUNG (v03-Regel): 6.8 = Beobachtbarkeit EINER Ausführung;
% 6.9 = reproduzierbare Prüfung des STANDES — kein zweites Mal
% Tests/Logs. Budget ~0,75 S. KEIN Listing (2/3–4 verbraucht).
% Zahlen-Politik: keine Zählungen (Testzahl -> 6.9, commit-
% gebunden).
% ============================================================

\section{Observability und technische Verifikation}
\label{sec:realisierung-observability}

Die Beobachtbarkeit aus
Abschnitt~\ref{sec:evidenzgebundene-verarbeitung} stützt sich auf
drei technische Ebenen: gespeicherte Laufartefakte, Protokolle der
Modell- und Werkzeugaufrufe sowie Telemetrie. Sie erfassen
unterschiedliche Ausschnitte der Ausführung und ergänzen sich bei
deren Rekonstruktion.

\emph{Laufartefakte.} Jeder Lauf erhält ein eigenes Verzeichnis
mit einem fortlaufenden Ereignisprotokoll und den anfallenden
Stufenartefakten. Durchlaufene Human Gates ergänzen Anfrage- und
Antwortartefakte; Prüf-Reparatur-Schleifen halten ihre Versuche
fest. Der Gesamtworkflow behandelt die Fehlerereignisse seines
Ereignisstroms ausdrücklich und weist den Lauf in diesen Fällen
als gescheitert aus. Die gespeicherten Belege erlauben es,
protokollierte Schritte nachträglich zu rekonstruieren, etwa Pause
und Wiederaufnahme aus
Abschnitt~\ref{sec:realisierung-durabilitaet}.

Diese Rekonstruktion reicht nur so weit wie die erhaltenen
Artefakte. Fehlende Belege für Konfigurationswechsel oder
Rücksetzungen in historischen Läufen bleiben deshalb
Nachweislücken, die Kapitel~\ref{chap:evaluation} ausweist.

\emph{Modell- und Werkzeugprotokolle.} Die gemeinsame
Chat-Pipeline protokolliert Modellaufrufe und Angaben zu ihrem
Eingabekontext. Diese Kontextanalyse erfasst die Struktur der
übergebenen Nachrichten und darin enthaltene Werkzeugaufrufe;
sie weist keine tatsächliche Nutzung oder Gewichtung der Inhalte
durch das Modell nach. Im detaillierten Protokollierungsmodus erfasst
die Chat-Pipeline auch die einzelnen Modellrunden innerhalb der
Werkzeugschleife;
im anderen Modus wird der Chat-Aufruf als Ganzes aufgezeichnet.
Die Werkzeugprotokolle erfassen Beginn, Rückgabe und Fehler eines
Aufrufs gesondert; eine erfolgreiche Rückgabe belegt für sich noch
keine beabsichtigte Zustandsänderung. Damit lässt sich die in
Abschnitt~\ref{sec:realisierung-tools-mcp} erklärte Trennung von
Ankündigung und Ausführung für die protokollierten Aufrufe prüfen.
Für ausgewählte agentische Knoten kann zusätzlich ein explizit
ausgegebener Analyse- oder Begründungstext des Modells
mitprotokolliert werden. Er dient der nachträglichen Analyse und
steuert den Ablauf nicht.

\emph{Telemetrie.} Die Chat-Pipeline nutzt die
OpenTelemetry-Anbindung des Frameworks. Sie erfasst Aufrufe und
Verbrauchsattribute nach den GenAI-Konventionen, also gemeinsamen
Bezeichnungen für Telemetriedaten generativer KI \cite{msmaf2026agentobservability}.
Bei aktiviertem Telemetrieexport schreiben projektspezifische
Exporter die Ablaufspuren (Traces) und Metriken in Dateien des
jeweiligen Laufs. Die Exporter übernehmen damit die lokale Ablage
der vom Framework bereitgestellten Signale. Eine abschließende
Aggregation führt vorhandene Telemetriedaten und Stufenartefakte
in einer Metrikdatei zusammen. Die erfassten Daten bilden eine
technische Grundlage für Kapitel~\ref{chap:evaluation}; ihre
Aussagekraft hängt von der gewählten Aufzeichnung und den
tatsächlich erhaltenen Belegen ab.

Wie automatisierte Tests diese Ausführungsbelege ergänzen und
welche Reichweite ihre Ergebnisse besitzen, erläutert
Abschnitt~\ref{sec:realisierung-reproduzierbarkeit-grenzen}.

% M18, 12.09.2026: Sprachliche Neufassung; historische Prüfnotizen
% aus der Vorfassung folgen als lokale Herkunftshinweise.
% [Rückverweis-Satz (Form §7; A1-Linie über 5.3). Leitfrage =
%  Gesamtvertrags-Endform (6.8/6.9-Abgrenzung v03).]
% [v02 (Kollegen-Punkte 1+2): ① „jeder Lauf erhält Gate-Artefakte/
%  Versuchshistorien" war zu universal (nicht jeder Lauf durchläuft
%  Gates/Repair — z. B. LLM-freie Läufe) ② „Lauf rekonstruierbar" ->
%  „PROTOKOLLIERTE Ausführungsschritte" (A1-Grenze: Beobachtbarkeit
%  macht feststellbar, nicht vollständig — Reasoning ist optional).
%  Merker: „Abschnitt 6.5" -> \ref GESETZT (Label-Pass 03.09.).]
% [runs/<stufe>/<runId>/ mit events.jsonl, gate/*.json
%  (GateAttempt-Historien = W2-Messdraht), logs/, decision-log;
%  Fehler-Events: R-18 (ExecutorFailedEvent/WorkflowErrorEvent
%  behandelt — „sonst endet der Lauf still", CLAUDE.md-Falle).
%  Selbst-Beleg: die V5-Forensik (§4d) rekonstruierte Pause/Resume/
%  Abschluss vollständig aus events.jsonl+checkpoints/ — der Satz
%  ist wörtlich wahr. Merker: „Abschnitt 6.5" -> \ref GESETZT (Label-Pass 03.09.).]
% [Freeze-Pass 04.09. (Kollegen-P1-10): „Messdraht" =
%  Entwicklerjargon; „als verlässliche Quelle" gestrichen
%  (Wertungs-Attribut, die Trennung selbst trägt die Aussage).]
% [v02 (Kollegen-Punkt 3): „JEDER Modellaufruf" scope-begrenzt auf
%  die gemeinsame Pipeline — ob der Steward-Dialog denselben
%  Messdraht nimmt, ist NICHT verifiziert (nicht behaupten).
%  Reasoning: „vom Modell explizit ausgegebener Text" (Mechanik =
%  Prompt-Appendix+Schema-Feld) — kein privilegierter Einblick in
%  den internen Schlussfolgerungsprozess impliziert.]
% [W1b „jeder Call am Messdraht": AgentChatPipelineBuilder.cs:44/73
%  (zwei Varianten); ChatDecisionLoggerMiddleware (ModelRound) ·
%  InputContextLoggerMiddleware · ToolCallLoggerMiddleware
%  (Code-Kommentar „die verlässliche Quelle für tatsächlich
%  ausgeführte Tools", ChatDecisionLoggerMiddleware.cs:51).
%  Reasoning: W1a captureReasoning-Schalter; Ledger „reasoning
%  log-only via Prompt-Appendix + Schema-Feld"
%  (SemanticLedgerExtractor.cs:194) — „log-only" => keine
%  Steuerungsfunktion, defensiv formuliert.
%  Merker: „Abschnitt 6.3" -> \ref GESETZT (Label-Pass 03.09.).]
% [„Kapitel~7" bewusst KLARTEXT: das Kapitel existiert noch nicht,
%  ein \ref würde als „??" rendern. Merker: bei Kap.-7-Anlage dort
%  \label{chap:evaluation} (o. ä.) setzen und HIER auf \ref
%  umstellen — gleicher Posten in 6.9.]
% [v02 (Kollegen-Punkt 4, WICHTIG): Die Abwesenheits-Behauptung
%  „Framework sieht keine laufbezogene Ablage vor" stützte sich nur
%  auf einen Klassen-Kommentar im EIGENEN Code — extern schwer
%  beweisbar und unnötig. Positive Fassung trägt dieselbe
%  MAF-vs-Eigenanteil-Aussage (MAF = OTel-Instrumentierung;
%  Projekt = Run-Senken + Aggregation).]
% [V6 #13 BESTÄTIGT: .UseOpenTelemetry(sourceName) + „OpenTelemetry
%  GenAI Semantic Conventions", gen_ai.client.token.usage etc.;
%  wir aktivieren auf CHAT-Ebene (Doppel-Erfassungs-Warnung der
%  Doku beachtet). Eigenanteil (Matrix B, bewusst custom):
%  OtelRunExporters (otel-traces.jsonl je Lauf — „MAF/OTel liefern
%  keine solche Senke", Klassen-Kommentar) + MetricsFinalizer ->
%  metrics.json. KEINE Wirksamkeits-/Messwert-Aussagen hier
%  (implementiert ≠ wirksam; Werte -> Kap. 7).]
% [v02 (Kollegen-Schlussfix): „Ausführungen belegbar" pauschal ->
%  epistemisch exakt (protokollierte Schritte, technische Belege).]
% [Schlussklammer = exakt die 6.8/6.9-Leitfragen-Abgrenzung
%  (v03-Regel: kein zweites Mal Tests/Logs); Übergang nummernfrei.]
